\documentclass[11pt,a4paper]{article}

% =========================================================
% COMPILATION
% =========================================================
% Compile using XeLaTeX or LuaLaTeX.
%
% Recommended:
%   xelatex voice-workflow-platform.tex
%   xelatex voice-workflow-platform.tex
%
% The second compilation resolves page references.

% =========================================================
% PACKAGES
% =========================================================
\usepackage[a4paper,margin=20mm,headheight=15pt]{geometry}
\usepackage{fontspec}
\usepackage{polyglossia}
\usepackage{microtype}
\usepackage{xcolor}
\usepackage{graphicx}
\usepackage{booktabs}
\usepackage{tabularx}
\usepackage{array}
\usepackage{array}
\usepackage{longtable}
\usepackage{enumitem}
\usepackage{titlesec}
\usepackage{fancyhdr}
\usepackage{lastpage}
\usepackage{hyperref}
\usepackage{listings}
\usepackage[most]{tcolorbox}
\usepackage{tikz}
\usepackage{float}
\usepackage{ragged2e}

\usetikzlibrary{
arrows.meta,
positioning,
fit,
calc,
shapes.geometric,
backgrounds
}

% =========================================================
% LANGUAGES AND FONTS
% =========================================================
\setdefaultlanguage{english}
\setotherlanguage{arabic}

\IfFontExistsTF{TeX Gyre Pagella}{
\setmainfont{TeX Gyre Pagella}
}{
\setmainfont{Times New Roman}
}

\IfFontExistsTF{TeX Gyre Heros}{
\setsansfont{TeX Gyre Heros}
}{
\setsansfont{Arial}
}

\IfFontExistsTF{Latin Modern Mono}{
\setmonofont{Latin Modern Mono}
}{
\setmonofont{Courier New}
}

% Arabic font fallback strategy.
\IfFontExistsTF{IBM Plex Sans Arabic}{
\newfontfamily\arabicfont[
Script=Arabic,
Scale=1.05
]{IBM Plex Sans Arabic}
}{
\IfFontExistsTF{Noto Sans Arabic}{
\newfontfamily\arabicfont[
Script=Arabic,
Scale=1.05
]{Noto Sans Arabic}
}{
\IfFontExistsTF{Amiri}{
\newfontfamily\arabicfont[
Script=Arabic,
Scale=1.05
]{Amiri}
}{
\newfontfamily\arabicfont[
Script=Arabic,
Scale=1.05
]{Arial}
}
}
}

% =========================================================
% DOCUMENT METADATA
% =========================================================
\newcommand{\ProjectName}{Enterprise Real-Time Voice-to-Workflow Platform}
\newcommand{\ProjectShortName}{Voice Workflow Runtime}
\newcommand{\DocumentVersion}{1.0}
\newcommand{\DocumentDate}{\today}
\newcommand{\PreparedFor}{Product, Engineering, AI, Security, and Operations Teams}

% =========================================================
% COLORS
% =========================================================
\definecolor{Primary}{HTML}{17324D}
\definecolor{PrimaryLight}{HTML}{EAF0F5}
\definecolor{Secondary}{HTML}{2D6A6A}
\definecolor{Accent}{HTML}{B88927}
\definecolor{Success}{HTML}{2E7D5B}
\definecolor{Warning}{HTML}{A66700}
\definecolor{Danger}{HTML}{A93A3A}
\definecolor{TextDark}{HTML}{17212B}
\definecolor{TextMuted}{HTML}{5F6B76}
\definecolor{Border}{HTML}{D7DEE5}
\definecolor{Surface}{HTML}{F7F9FB}
\definecolor{CodeBackground}{HTML}{F4F6F8}

% =========================================================
% PDF AND LINKS
% =========================================================
\hypersetup{
pdftitle={\ProjectName},
pdfauthor={Project Team},
pdfsubject={Product and Technical Architecture Specification},
colorlinks=true,
linkcolor=Primary,
urlcolor=Secondary,
citecolor=Primary
}

% =========================================================
% HEADER AND FOOTER
% =========================================================
\pagestyle{fancy}
\fancyhf{}
\fancyhead[L]{\sffamily\small\ProjectShortName}
\fancyhead[R]{\sffamily\small Version \DocumentVersion}
\fancyfoot[L]{\sffamily\footnotesize Confidential --- Internal Use}
\fancyfoot[R]{\sffamily\footnotesize Page \thepage\ of \pageref{LastPage}}

\renewcommand{\headrulewidth}{0.4pt}
\renewcommand{\footrulewidth}{0.4pt}

% =========================================================
% SECTION STYLING
% =========================================================
\titleformat{\section}
{\Large\sffamily\bfseries\color{Primary}}
{\thesection}
{0.7em}
{}

\titleformat{\subsection}
{\large\sffamily\bfseries\color{Primary}}
{\thesubsection}
{0.7em}
{}

\titleformat{\subsubsection}
{\normalsize\sffamily\bfseries\color{Secondary}}
{\thesubsubsection}
{0.7em}
{}

\titlespacing*{\section}{0pt}{22pt}{10pt}
\titlespacing*{\subsection}{0pt}{17pt}{7pt}
\titlespacing*{\subsubsection}{0pt}{13pt}{5pt}

% =========================================================
% TABLE HELPERS
% =========================================================
\newcolumntype{Y}{>{\RaggedRight\arraybackslash}X}
\newcolumntype{L}[1]{>{\RaggedRight\arraybackslash}p{#1}}
\renewcommand{\arraystretch}{1.25}

% =========================================================
% LIST STYLING
% =========================================================
\setlist[itemize]{
leftmargin=1.5em,
itemsep=3pt,
topsep=4pt
}

\setlist[enumerate]{
leftmargin=1.7em,
itemsep=3pt,
topsep=4pt
}

% =========================================================
% CODE LISTINGS
% =========================================================
\lstdefinestyle{projectcode}{
backgroundcolor=\color{CodeBackground},
basicstyle=\ttfamily\footnotesize,
keywordstyle=\bfseries\color{Primary},
stringstyle=\color{Success},
commentstyle=\itshape\color{TextMuted},
numberstyle=\tiny\color{TextMuted},
numbers=left,
numbersep=8pt,
frame=single,
rulecolor=\color{Border},
breaklines=true,
breakatwhitespace=false,
showstringspaces=false,
keepspaces=true,
columns=fullflexible,
tabsize=2,
xleftmargin=3pt,
xrightmargin=3pt,
aboveskip=8pt,
belowskip=8pt
}

\lstset{style=projectcode}

% =========================================================
% BOXES
% =========================================================
\newtcolorbox{executivebox}{
enhanced,
breakable,
colback=PrimaryLight,
colframe=Primary,
boxrule=0.8pt,
arc=2mm,
left=4mm,
right=4mm,
top=3mm,
bottom=3mm,
fonttitle=\sffamily\bfseries,
title=Executive Summary
}

\newtcolorbox{decisionbox}[1][]{
enhanced,
breakable,
colback=Surface,
colframe=Secondary,
boxrule=0.7pt,
arc=2mm,
left=4mm,
right=4mm,
top=3mm,
bottom=3mm,
fonttitle=\sffamily\bfseries,
title={#1}
}

\newtcolorbox{warningbox}[1][]{
enhanced,
breakable,
colback=Warning!8,
colframe=Warning,
boxrule=0.7pt,
arc=2mm,
left=4mm,
right=4mm,
top=3mm,
bottom=3mm,
fonttitle=\sffamily\bfseries,
title={#1}
}

\newtcolorbox{securitybox}[1][]{
enhanced,
breakable,
colback=Danger!5,
colframe=Danger,
boxrule=0.7pt,
arc=2mm,
left=4mm,
right=4mm,
top=3mm,
bottom=3mm,
fonttitle=\sffamily\bfseries,
title={#1}
}

% =========================================================
% TIKZ STYLES
% =========================================================
\tikzset{
component/.style={
rectangle,
rounded corners=2mm,
minimum width=30mm,
minimum height=11mm,
align=center,
draw=Primary,
fill=PrimaryLight,
font=\sffamily\small
},
service/.style={
rectangle,
rounded corners=2mm,
minimum width=31mm,
minimum height=12mm,
align=center,
draw=Secondary,
fill=Secondary!8,
font=\sffamily\small
},
storage/.style={
cylinder,
shape border rotate=90,
aspect=0.22,
minimum width=23mm,
minimum height=13mm,
align=center,
draw=Accent,
fill=Accent!10,
font=\sffamily\small
},
flow/.style={
-{Latex[length=2.2mm]},
thick,
draw=TextMuted
},
asyncflow/.style={
-{Latex[length=2.2mm]},
thick,
dashed,
draw=TextMuted
}
}

% =========================================================
% DOCUMENT
% =========================================================
\begin{document}

% =========================================================
% TITLE PAGE
% =========================================================
\begin{titlepage}
\thispagestyle{empty}

\vspace*{16mm}

{\sffamily\Large\color{Secondary}
PRODUCT AND TECHNICAL ARCHITECTURE
}

\vspace{8mm}

{\sffamily\bfseries\fontsize{27}{33}\selectfont\color{Primary}
\ProjectName
}

\vspace{5mm}

{\Large\color{TextMuted}
A configurable enterprise platform that converts live Arabic speech
into validated, secure, executable business workflows.
}

\vspace{14mm}

\begin{tcolorbox}[
colback=Surface,
colframe=Border,
boxrule=0.6pt,
arc=2mm,
left=5mm,
right=5mm,
top=4mm,
bottom=4mm
]
\begin{tabularx}{\textwidth}{L{37mm}Y}
\textbf{Document Version} & \DocumentVersion \\
\textbf{Prepared For} & \PreparedFor \\
\textbf{Document Date} & \DocumentDate \\
\textbf{Status} & Initial architecture and product definition \\
\textbf{Classification} & Confidential --- Internal Use \\
\end{tabularx}
\end{tcolorbox}

\vfill

\begin{center}
{\sffamily\small\color{TextMuted}
This document defines the product vision, system boundaries,
real-time architecture, security model, API contracts,
delivery roadmap, and acceptance criteria.
}
\end{center}
\end{titlepage}

\tableofcontents
\newpage

% =========================================================
\section{Executive Summary}

\begin{executivebox}
\ProjectName{} is a multi-tenant enterprise platform that allows each
organization to define its own business workflows, data schemas,
terminology, validation rules, permissions, and integration targets.

An employee can speak naturally in Arabic or Egyptian Arabic while the
platform processes the speech as a real-time stream. The system produces
a continuously updated structured workflow draft, validates it against
the organization's configuration, asks for missing information when
necessary, and executes the final action only after the configured
confirmation and approval requirements are satisfied.

The platform is not only a speech-to-text service and is not a generic
large-language-model wrapper. Its primary value is the combination of:

\begin{itemize}
\item Real-time streaming speech processing.
\item Tenant-defined JSON schemas and workflow definitions.
\item Incremental structured-data extraction.
\item Deterministic business validation.
\item Enterprise identity and permission enforcement.
\item Secure workflow execution.
\item Webhooks, APIs, queues, and n8n integrations.
\item Auditing, replay, observability, and quality evaluation.
\end{itemize}
\end{executivebox}

% =========================================================
\section{Problem Definition}

Operational employees frequently communicate business events using voice
because typing structured forms is slow, inconvenient, or impractical in
field environments.

Examples include:

\begin{itemize}
\item Drivers recording sales, purchases, collections, or expenses.
\item Warehouse employees recording stock movements.
\item Sales representatives recording customer orders.
\item Maintenance teams recording incidents and completed work.
\item Recruitment teams updating candidate or worker states.
\item Healthcare and service teams recording operational notes.
\end{itemize}

A typical spoken statement may contain multiple entities, calculations,
corrections, and business decisions.

For example:

\begin{quote}
\textarabic{
اشتريت من شركة النور ثلاثة طن أرز، الطن بأربعة وعشرين ألف،
ودفعت ثلاثين ألف كاش، والباقي آجل، والنقل ألف ومائتين جنيه.
}
\end{quote}

Traditional speech-to-text systems return plain text. They do not reliably
determine:

\begin{itemize}
\item Which business workflow the employee intends to execute.
\item Which fields belong to the workflow.
\item Whether the values comply with company rules.
\item Whether referenced suppliers, customers, products, or accounts exist.
\item Whether the employee is authorized to perform the action.
\item Whether confirmation or approval is required.
\item Which external application should receive the final result.
\end{itemize}

% =========================================================
\section{Product Vision}

\subsection{Vision Statement}

\begin{decisionbox}[Product Vision]
Provide enterprises with a configurable runtime that converts natural
spoken language into validated business actions without requiring each
organization to build its own speech, AI, workflow, security, and
integration infrastructure.
\end{decisionbox}

\subsection{Product Positioning}

The platform should be positioned as:

\begin{quote}
\textbf{Enterprise Real-Time Voice Workflow Infrastructure}
\end{quote}

It should not be positioned as:

\begin{itemize}
\item A generic transcription tool.
\item A chatbot builder.
\item A fixed accounting application.
\item A general-purpose automation platform.
\item An unvalidated AI action executor.
\end{itemize}

\subsection{Core Differentiation}

The platform's competitive advantage is the combination of:

\begin{enumerate}
\item Streaming voice processing with low perceived latency.
\item Configurable enterprise workflows.
\item Constrained JSON output.
\item Deterministic validation and entity resolution.
\item Secure and auditable execution.
\item Arabic and Egyptian Arabic operational language support.
\end{enumerate}

% =========================================================
\section{Target Customers and Use Cases}

\subsection{Primary Customer Segments}

\begin{tabularx}{\textwidth}{L{36mm}Y Y}
\toprule
\textbf{Segment} & \textbf{Operational Need} & \textbf{Example Workflows} \\
\midrule
Distribution companies
& Fast field data capture
& Sales, collections, expenses, stock transfer \\

Agricultural companies
& Field and warehouse operations
& Purchase, delivery, batch movement, inspection \\

Logistics companies
& Driver and shipment reporting
& Delivery, incident, fuel expense, cash collection \\

ERP and accounting vendors
& Voice-enabled transaction entry
& Purchase, sale, payment, receipt, journal draft \\

Maintenance companies
& Hands-free field reporting
& Work order, incident, parts used, completion report \\

Recruitment agencies
& Worker and candidate workflow updates
& Candidate status, document collection, assignment \\

Healthcare operations
& Structured operational recording
& Visit notes, inventory usage, service completion \\
\bottomrule
\end{tabularx}

\subsection{Primary Personas}

\begin{itemize}
\item \textbf{Enterprise Administrator:}
defines workflows, schemas, rules, integrations, and access policies.

\item \textbf{Operational Employee:}
speaks naturally and reviews the generated structured action.

\item \textbf{Supervisor or Approver:}
reviews high-risk actions before execution.

\item \textbf{Integration Engineer:}
connects the platform to internal systems, APIs, and n8n.

\item \textbf{Security Administrator:}
manages identity, keys, retention, access controls, and audit policies.

\item \textbf{Quality Analyst:}
reviews extraction quality, failed sessions, corrections, and model metrics.
\end{itemize}

% =========================================================
\section{Core Product Concepts}

\subsection{Tenant}

A tenant represents one customer organization. All workflow definitions,
data, credentials, logs, terminology, and integrations are isolated by
tenant.

\subsection{Workflow Definition}

A workflow definition describes a business operation such as:

\begin{itemize}
\item Purchase
\item Sale
\item Income
\item Expense
\item Receipt
\item Payment
\item Stock transfer
\item Maintenance request
\item Candidate status update
\end{itemize}

Each workflow contains:

\begin{itemize}
\item Unique key and display name.
\item Business description.
\item JSON Schema.
\item Required and optional fields.
\item Trigger phrases and terminology.
\item Examples.
\item Validation rules.
\item Entity-resolution rules.
\item Permission requirements.
\item Execution policy.
\item Integration destination.
\item Version.
\end{itemize}

\subsection{Voice Session}

A voice session represents one continuous employee interaction. It stores:

\begin{itemize}
\item Tenant and user identity.
\item Audio stream metadata.
\item Transcript revisions.
\item Workflow candidates.
\item Current structured draft.
\item Missing fields.
\item Validation results.
\item Confirmation state.
\item Final execution result.
\end{itemize}

\subsection{Workflow Draft}

A workflow draft is an incomplete or complete structured representation
of the employee's spoken intent.

It is mutable while the employee is speaking and immutable after final
commit.

\subsection{Execution Policy}

The execution policy controls what must happen before the platform sends
or executes the final workflow.

Supported modes should include:

\begin{itemize}
\item Automatic execution.
\item User confirmation.
\item PIN or biometric confirmation.
\item Supervisor approval.
\item Dual approval.
\item Draft creation only.
\end{itemize}

% =========================================================
\section{End-to-End User Experience}

\subsection{Administrator Experience}

The enterprise administrator opens the dashboard and performs the
following process:

\begin{enumerate}
\item Creates a new workflow.
\item Defines its name, description, and business purpose.
\item Creates or imports its JSON Schema.
\item Marks required and optional fields.
\item Adds example phrases.
\item Defines company-specific synonyms.
\item Adds validation rules.
\item Configures database or API entity resolution.
\item Assigns roles and permissions.
\item Selects an execution policy.
\item Connects a webhook, API, queue, or n8n workflow.
\item Tests the workflow in a voice playground.
\item Publishes a version for production use.
\end{enumerate}

\subsection{Employee Experience}

The operational employee:

\begin{enumerate}
\item Opens the voice-enabled application.
\item Starts a secure voice session.
\item Speaks naturally.
\item Sees an immediate partial transcript.
\item Sees a structured draft update while speaking.
\item Corrects information naturally when necessary.
\item Provides missing information when asked.
\item Reviews the final structured result.
\item Confirms the action.
\item Receives a success, rejection, or approval-pending result.
\end{enumerate}

\subsection{Example Interaction}

The employee says:

\begin{quote}
\textarabic{
بعت للعميل أحمد ثلاثة طن أرز، الطن باثنين وعشرين ألف،
واستلمت منه عشرة آلاف كاش، والباقي آجل.
}
\end{quote}

The platform incrementally identifies:

\begin{itemize}
\item Workflow: Sale
\item Customer: Ahmed
\item Product: Rice
\item Quantity: 3
\item Unit: Ton
\item Unit price: 22,000
\item Paid amount: 10,000
\item Remaining amount: 56,000
\item Payment status: Partial
\end{itemize}

% =========================================================
\section{Workflow Lifecycle}

The workflow session follows a controlled lifecycle.

\begin{center}
\begin{tikzpicture}[
node distance=8mm and 6mm,
every node/.style={font=\sffamily\small}
]
\node[component] (listening) {LISTENING};
\node[component, right=of listening] (draft) {DRAFT};
\node[component, right=of draft] (ready) {READY};
\node[component, right=of ready] (confirmed) {CONFIRMED};
\node[component, right=of confirmed] (committed) {COMMITTED};

\draw[flow] (listening) -- (draft);
\draw[flow] (draft) -- (ready);
\draw[flow] (ready) -- (confirmed);
\draw[flow] (confirmed) -- (committed);

\node[component, below=12mm of ready, draw=Danger, fill=Danger!6] (rejected)
{REJECTED};

\draw[flow, draw=Danger] (draft) -- (rejected);
\draw[flow, draw=Danger] (ready) -- (rejected);
\draw[flow, draw=Danger] (confirmed) -- (rejected);
\end{tikzpicture}
\end{center}

\begin{tabularx}{\textwidth}{L{30mm}Y}
\toprule
\textbf{State} & \textbf{Meaning} \\
\midrule
LISTENING
& Audio is being received and transcribed. \\

DRAFT
& The platform has identified a possible workflow and partial fields. \\

READY
& Required information is complete and deterministic validation has passed. \\

CONFIRMED
& The user or approver has confirmed execution. \\

COMMITTED
& The final action has been delivered or executed successfully. \\

REJECTED
& The action was rejected due to validation, authorization, approval, or integration failure. \\
\bottomrule
\end{tabularx}

\begin{warningbox}[Critical Execution Rule]
Partial transcripts and draft updates must never directly trigger
irreversible business actions.

Only a confirmed and validated final workflow may reach the execution
dispatcher.
\end{warningbox}

% =========================================================
\section{Real-Time Processing Architecture}

\subsection{High-Level Architecture}

\begin{figure}[H]
\centering
\resizebox{\textwidth}{!}{
\begin{tikzpicture}[
node distance=10mm and 13mm
]
\node[component] (client) {Mobile or Web\\Client};

\node[service, right=of client] (gateway)
{Realtime\\Gateway};

\node[service, right=of gateway] (speech)
{Streaming Speech\\Service};

\node[service, right=of speech] (runtime)
{Workflow\\Runtime};

\node[service, right=of runtime] (validation)
{Validation and\\Entity Resolution};

\node[service, right=of validation] (dispatcher)
{Action\\Dispatcher};

\node[storage, below=19mm of gateway] (redis)
{Redis};

\node[storage, below=19mm of runtime] (postgres)
{PostgreSQL};

\node[service, below=19mm of validation] (queue)
{Event Queue};

\node[service, below=19mm of dispatcher] (webhook)
{Webhook and\\Integration Workers};

\draw[flow] (client) -- node[above, font=\scriptsize]{Audio / events} (gateway);
\draw[flow] (gateway) -- (speech);
\draw[flow] (speech) -- node[above, font=\scriptsize]{Partial transcript} (runtime);
\draw[flow] (runtime) -- (validation);
\draw[flow] (validation) -- node[above, font=\scriptsize]{Validated action} (dispatcher);

\draw[flow] (gateway) -- (redis);
\draw[flow] (runtime) -- (postgres);

\draw[asyncflow] (validation) -- (queue);
\draw[asyncflow] (dispatcher) -- (webhook);
\draw[asyncflow] (queue) -- (webhook);

\draw[flow, bend left=32] (runtime.north) to
node[above, font=\scriptsize]{Draft updates} (gateway.north);

\draw[flow, bend left=36] (gateway.north) to
node[above, font=\scriptsize]{Realtime UI events} (client.north);
\end{tikzpicture}
}
\caption{High-level real-time and asynchronous processing architecture.}
\end{figure}

\subsection{Realtime Gateway}

The Realtime Gateway is responsible for:

\begin{itemize}
\item WebSocket or WebRTC session management.
\item Authentication and tenant resolution.
\item Audio chunk sequencing.
\item Backpressure and connection health.
\item Reconnection and session recovery.
\item Rate limiting and abuse prevention.
\item Routing sessions to speech-processing workers.
\item Returning transcript and workflow events to the client.
\end{itemize}

The gateway should remain stateless where possible. Hot session state
should be stored in Redis.

\subsection{Streaming Speech Service}

The speech service is responsible for:

\begin{itemize}
\item Audio format validation.
\item Voice activity detection.
\item Noise and silence handling.
\item Streaming speech-to-text.
\item Partial transcript generation.
\item Final transcript generation.
\item Transcript stability scoring.
\item Arabic number normalization.
\item Egyptian Arabic terminology normalization.
\end{itemize}

\subsection{Workflow Runtime}

The Workflow Runtime is responsible for:

\begin{itemize}
\item Loading tenant workflow definitions.
\item Selecting workflow candidates.
\item Extracting structured fields.
\item Updating the current draft.
\item Handling corrections and revisions.
\item Tracking source spans and confidence.
\item Identifying missing required fields.
\item Deciding when deeper model reasoning is required.
\end{itemize}

\subsection{Validation and Entity Resolution}

This service provides deterministic validation after AI extraction.

It validates:

\begin{itemize}
\item JSON Schema compliance.
\item Required fields.
\item Numeric ranges.
\item Cross-field business rules.
\item Entity existence.
\item User permissions.
\item Tenant restrictions.
\item Duplicate-action risks.
\item Accounting and inventory constraints.
\end{itemize}

\subsection{Action Dispatcher}

The Action Dispatcher executes only finalized workflows.

Supported delivery mechanisms should include:

\begin{itemize}
\item Signed webhook delivery.
\item REST API invocation.
\item Message queue publication.
\item n8n webhook trigger.
\item Internal connector invocation.
\item Draft-only persistence.
\end{itemize}

% =========================================================
\section{Low-Latency Processing Strategy}

\subsection{Fast Layer and Reasoning Layer}

The system should not send every partial word to a large language model.

Instead, it should use two layers.

\subsubsection{Fast Layer}

The fast layer runs continuously and performs:

\begin{itemize}
\item Partial transcription.
\item Number normalization.
\item Keyword recognition.
\item Basic entity extraction.
\item Candidate workflow detection.
\item Session state updates.
\end{itemize}

\subsubsection{Reasoning Layer}

The reasoning layer runs at semantic checkpoints, including:

\begin{itemize}
\item A stable phrase or sentence boundary.
\item A short period of silence.
\item Detection of a new business entity.
\item Detection of a correction.
\item Completion of required workflow information.
\item Ambiguity requiring structured reasoning.
\item Final end-of-speech processing.
\end{itemize}

\subsection{Target Latency Budget}

\begin{tabularx}{\textwidth}{Y L{34mm} Y}
\toprule
\textbf{Stage} & \textbf{Target} & \textbf{User-Visible Result} \\
\midrule
Audio chunk upload
& 20--100 ms
& Continuous connection \\

Partial speech recognition
& 100--300 ms
& Live transcript \\

Fast entity extraction
& 10--60 ms
& Immediate field updates \\

Draft UI update
& Below 100 ms
& Responsive structured form \\

Structured model reasoning
& 200--900 ms
& Corrected workflow draft \\

Final deterministic validation
& 20--250 ms
& Ready or validation-error state \\
\bottomrule
\end{tabularx}

These figures are architectural targets and should be verified through
load testing against the selected speech and language models.

% =========================================================
\section{Workflow Configuration Model}

\subsection{Example Workflow Definition}

\begin{lstlisting}[language={},caption={Example tenant workflow definition}]
{
"key": "purchase",
"version": 3,
"name": "Purchase",
"description": "Record the purchase of goods from a supplier",
"status": "published",

"triggerPhrases": [
"I purchased",
"we bought",
"goods received from supplier"
],

"schema": {
"type": "object",
"required": ["supplierId", "items"],
"properties": {
"supplierId": {
"type": "string"
},
"items": {
"type": "array",
"minItems": 1,
"items": {
"type": "object",
"required": [
"productId",
"quantity",
"unitId",
"unitPrice"
],
"properties": {
"productId": {
"type": "string"
},
"quantity": {
"type": "number",
"exclusiveMinimum": 0
},
"unitId": {
"type": "string"
},
"unitPrice": {
"type": "number",
"minimum": 0
}
}
}
},
"paidAmount": {
"type": "number",
"minimum": 0
},
"paymentMethod": {
"type": "string",
"enum": ["cash", "bank", "credit"]
},
"expenses": {
"type": "array"
}
}
},

"rules": [
{
"key": "payment_not_above_total",
"expression": "paidAmount <= invoiceTotal",
"severity": "error"
}
],

"executionPolicy": {
"mode": "confirm",
"requiresPin": false,
"requiresSupervisorApproval": false
},

"integration": {
"type": "webhook",
"destinationId": "dest_01KEXAMPLE"
}
}
\end{lstlisting}

\subsection{Workflow Versioning}

Workflow definitions must be versioned.

Publishing a new version must not alter historical sessions created using
older versions.

Each session must store:

\begin{itemize}
\item Workflow definition ID.
\item Workflow version.
\item Model configuration version.
\item Validation-rules version.
\item Integration-destination version.
\end{itemize}

\subsection{Schema Builder}

The dashboard should provide both:

\begin{itemize}
\item A visual schema builder for non-technical administrators.
\item An advanced JSON Schema editor for engineers.
\end{itemize}

The visual builder should support:

\begin{itemize}
\item Text, number, boolean, date, enum, object, and array fields.
\item Required and optional fields.
\item Default values.
\item Numeric limits.
\item String patterns.
\item Field descriptions.
\item Entity-reference fields.
\item Conditional fields.
\item Repeatable line items.
\end{itemize}

% =========================================================
\section{Realtime Protocol}

\subsection{Connection Establishment}

The client creates a session using a standard authenticated API request,
then opens a realtime connection.

\begin{lstlisting}[language={},caption={Create a voice session}]
POST /v1/realtime/sessions

{
"locale": "ar-EG",
"allowedWorkflowKeys": [
"purchase",
"sale",
"expense",
"income"
],
"audio": {
"format": "pcm16",
"sampleRate": 16000,
"channels": 1
}
}
\end{lstlisting}

\subsection{Client Audio Event}

Binary frames are preferred for audio payloads to avoid Base64 overhead.

A corresponding metadata event may use the following structure:

\begin{lstlisting}[language={},caption={Audio chunk metadata event}]
{
"type": "audio.chunk",
"sequence": 142,
"durationMs": 100,
"timestamp": "2026-07-23T18:30:11.420Z"
}
\end{lstlisting}

\subsection{Partial Transcript Event}

\begin{lstlisting}[language={},caption={Partial transcript event}]
{
"type": "transcript.partial",
"sessionId": "vs_01KEXAMPLE",
"revision": 12,
"text": "I purchased three tons of rice",
"isStable": false,
"confidence": 0.82
}
\end{lstlisting}

\subsection{Workflow Draft Update}

\begin{lstlisting}[language={},caption={Incremental workflow draft update}]
{
"type": "workflow.draft.updated",
"sessionId": "vs_01KEXAMPLE",
"revision": 7,
"workflow": {
"key": "purchase",
"state": "selected",
"confidence": 0.91
},
"data": {
"supplierId": "supplier_102",
"items": [
{
"productId": "product_rice_01",
"quantity": 3,
"unitId": "ton"
}
]
},
"missingFields": [
{
"path": "items[0].unitPrice",
"reason": "required"
}
]
}
\end{lstlisting}

\subsection{Final Workflow Event}

\begin{lstlisting}[language={},caption={Workflow ready for confirmation}]
{
"type": "workflow.ready",
"sessionId": "vs_01KEXAMPLE",
"revision": 14,
"workflowKey": "purchase",
"data": {
"supplierId": "supplier_102",
"items": [
{
"productId": "product_rice_01",
"quantity": 3,
"unitId": "ton",
"unitPrice": 24000,
"lineTotal": 72000
}
],
"paidAmount": 30000,
"remainingAmount": 42000,
"paymentMethod": "cash"
},
"validation": {
"valid": true,
"errors": [],
"warnings": []
},
"confirmationRequired": true
}
\end{lstlisting}

% =========================================================
\section{Corrections and Transcript Revisions}

Employees frequently correct themselves while speaking.

Example:

\begin{quote}
\textarabic{
بعت خمسة طن... لا، ثلاثة طن.
}
\end{quote}

The system must support revision-based updates rather than append-only
transcription.

\begin{lstlisting}[language={},caption={Transcript correction event}]
{
"type": "transcript.revision",
"revision": 18,
"replacement": {
"startToken": 4,
"endToken": 6,
"oldText": "five tons",
"newText": "three tons"
}
}
\end{lstlisting}

The Workflow Runtime should then update the affected field.

\begin{lstlisting}[language={},caption={Structured field correction}]
{
"type": "workflow.field.corrected",
"revision": 9,
"path": "items[0].quantity",
"oldValue": 5,
"newValue": 3,
"reason": "user_correction",
"confidence": 0.96
}
\end{lstlisting}

Each extracted field should retain provenance metadata:

\begin{itemize}
\item Transcript source range.
\item Transcript revision.
\item Extraction method.
\item Confidence.
\item Last updated timestamp.
\item Whether the value was user-confirmed.
\end{itemize}

% =========================================================
\section{Entity Resolution}

The platform should not assume that extracted names represent valid
business entities.

For example, ``Ahmed'' may match multiple customers.

Entity resolution should support:

\begin{itemize}
\item Exact matching.
\item Alias matching.
\item Phonetic matching.
\item Fuzzy text matching.
\item Recent-entity prioritization.
\item User-assigned entity scope.
\item External API lookup.
\item Human clarification.
\end{itemize}

When multiple entities are possible, the system should return candidates
instead of silently selecting one.

\begin{lstlisting}[language={},caption={Entity ambiguity event}]
{
"type": "workflow.clarification.required",
"field": "customerId",
"question": "Which Ahmed customer do you mean?",
"candidates": [
{
"id": "customer_11",
"label": "Ahmed Hassan - Cairo",
"confidence": 0.78
},
{
"id": "customer_42",
"label": "Ahmed Ali - Tanta",
"confidence": 0.72
}
]
}
\end{lstlisting}

% =========================================================
\section{Validation Model}

\subsection{Validation Layers}

Validation should occur through multiple layers.

\begin{enumerate}
\item \textbf{Schema Validation}

Confirms that the output structure and field types match the published
JSON Schema.

\item \textbf{Business Rule Validation}

Evaluates rules such as payment limits, quantity limits, account state,
and workflow-specific constraints.

\item \textbf{Entity Validation}

Confirms that referenced customers, suppliers, products, warehouses,
and accounts exist and are active.

\item \textbf{Authorization Validation}

Confirms that the authenticated employee may create, confirm, approve,
or execute the workflow.

\item \textbf{Risk Validation}

Detects duplicate actions, unusual values, policy violations, or actions
requiring additional approval.

\item \textbf{Integration Validation}

Confirms that the destination configuration is active and compatible
with the final payload version.
\end{enumerate}

\subsection{Validation Error Structure}

\begin{lstlisting}[language={},caption={Structured validation error}]
{
"valid": false,
"errors": [
{
"code": "PAID_AMOUNT_EXCEEDS_TOTAL",
"path": "paidAmount",
"message": "Paid amount cannot exceed the invoice total",
"severity": "error",
"source": "business_rule"
}
],
"warnings": []
}
\end{lstlisting}

% =========================================================
\section{Security Architecture}

\begin{securitybox}[Security Principle]
The language model must never be treated as an authorization system,
validation authority, or trusted execution engine.
\end{securitybox}

\subsection{Identity and Access Management}

The platform should support:

\begin{itemize}
\item OpenID Connect and SAML for enterprise identity.
\item JWT access tokens with short expiration.
\item Refresh-token rotation.
\item Multi-factor authentication.
\item Role-based access control.
\item Attribute-based policies for sensitive operations.
\item Device registration and revocation.
\item Service accounts for integrations.
\end{itemize}

\subsection{Tenant Isolation}

Tenant isolation must be enforced at:

\begin{itemize}
\item API authorization.
\item Database queries.
\item Redis keys.
\item Queue messages.
\item Object-storage paths.
\item Encryption context.
\item Metrics and logs.
\item Model prompts and context.
\end{itemize}

Every internal request should include a verified tenant context. Tenant
identifiers received from clients must not be trusted without
authentication and authorization.

\subsection{Data Protection}

The system should implement:

\begin{itemize}
\item TLS for all external and internal traffic.
\item Encryption at rest.
\item Tenant-specific encryption keys for enterprise plans.
\item Secret management outside application code.
\item Configurable audio-retention policies.
\item Transcript redaction.
\item Personally identifiable information controls.
\item Secure deletion workflows.
\end{itemize}

\subsection{Webhook Security}

Webhook delivery should support:

\begin{itemize}
\item HMAC signatures.
\item Timestamped requests.
\item Replay-attack protection.
\item Secret rotation.
\item IP allowlists.
\item Mutual TLS for enterprise customers.
\item Idempotency keys.
\item Retry and dead-letter policies.
\item Delivery logs.
\end{itemize}

\subsection{Prompt and Model Security}

The system must protect against:

\begin{itemize}
\item Prompt injection in spoken content.
\item Attempts to override workflow constraints.
\item Unauthorized workflow selection.
\item Data exfiltration through model responses.
\item Cross-tenant context leakage.
\item Malicious entity names or external content.
\end{itemize}

The model should receive only:

\begin{itemize}
\item The workflows allowed for the current user.
\item The minimum required tenant terminology.
\item The current transcript window.
\item The current structured draft.
\item Explicit system-controlled output constraints.
\end{itemize}

% =========================================================
\section{Auditability and Compliance}

Every session should produce a complete audit record containing:

\begin{itemize}
\item Tenant ID.
\item User ID.
\item Device ID.
\item Session timestamps.
\item Workflow definition and version.
\item Transcript revisions.
\item Draft revisions.
\item Model and prompt versions.
\item Validation results.
\item Confirmation evidence.
\item Approval decisions.
\item Final payload.
\item Integration-delivery attempts.
\item Final execution outcome.
\end{itemize}

Audio storage must be configurable per tenant:

\begin{itemize}
\item Do not store.
\item Store temporarily.
\item Store only failed sessions.
\item Store with explicit employee consent.
\item Store according to customer-defined retention periods.
\end{itemize}

% =========================================================
\section{Scalability and Availability}

\subsection{Independent Scaling Units}

The following components should scale independently:

\begin{itemize}
\item Realtime Gateway: concurrent connections.
\item Speech Workers: concurrent audio streams.
\item Workflow Runtime: extraction and reasoning requests.
\item Validation Workers: database and external lookups.
\item Webhook Workers: delivery throughput.
\item Audit Pipeline: event-ingestion volume.
\end{itemize}

\subsection{Hot and Durable State}

\begin{tabularx}{\textwidth}{L{35mm}Y Y}
\toprule
\textbf{Storage} & \textbf{Purpose} & \textbf{Examples} \\
\midrule
Redis
& Short-lived realtime state
& Session state, connection routing, rate limits \\

PostgreSQL
& Durable relational state
& Tenants, workflows, sessions, approvals, audit metadata \\

Object storage
& Large binary objects
& Audio, attachments, exported audit packages \\

Event queue
& Durable asynchronous processing
& Webhook jobs, analytics, quality evaluation \\
\bottomrule
\end{tabularx}

\subsection{Failure Handling}

The architecture should support:

\begin{itemize}
\item Realtime reconnection using session resume tokens.
\item Audio sequence-number verification.
\item Duplicate chunk detection.
\item Idempotent finalization.
\item At-least-once asynchronous delivery.
\item Idempotency keys for downstream execution.
\item Dead-letter queues.
\item Controlled retry with exponential backoff.
\item Circuit breakers for external dependencies.
\end{itemize}

% =========================================================
\section{Suggested Data Model}

\begin{longtable}{L{38mm}L{46mm}L{76mm}}
\toprule
\textbf{Entity} & \textbf{Purpose} & \textbf{Important Fields} \\
\midrule
\endhead

Tenant
& Customer organization
& ID, name, plan, region, security policy \\

User
& Human platform user
& Tenant, identity provider, roles, status \\

Device
& Registered client device
& User, fingerprint, status, last seen \\

WorkflowDefinition
& Logical workflow
& Tenant, key, name, current published version \\

WorkflowVersion
& Immutable workflow version
& Schema, rules, examples, execution policy \\

TerminologyEntry
& Tenant vocabulary
& Phrase, normalized value, entity type \\

IntegrationDestination
& External action destination
& Type, endpoint, secret reference, status \\

VoiceSession
& Realtime interaction
& User, workflow candidate, status, timestamps \\

TranscriptRevision
& Transcript history
& Session, revision, text, stability, confidence \\

WorkflowDraftRevision
& Structured draft history
& Session, revision, payload, missing fields \\

ValidationResult
& Validation outcome
& Session, rule, severity, field path \\

Confirmation
& User confirmation evidence
& User, method, timestamp, device \\

Approval
& Supervisor decision
& Approver, status, reason, timestamp \\

ExecutionAttempt
& Delivery or execution attempt
& Destination, attempt number, response, status \\

AuditEvent
& Immutable operational audit
& Actor, action, target, timestamp, metadata \\
\bottomrule
\end{longtable}

% =========================================================
\section{API Surface}

\subsection{Management APIs}

\begin{itemize}
\item \texttt{POST /v1/workflows}
\item \texttt{GET /v1/workflows}
\item \texttt{GET /v1/workflows/{workflowId}}
\item \texttt{POST /v1/workflows/{workflowId}/versions}
\item \texttt{POST /v1/workflows/{workflowId}/publish}
\item \texttt{POST /v1/workflows/{workflowId}/test}
\item \texttt{POST /v1/integrations}
\item \texttt{POST /v1/terminology}
\item \texttt{GET /v1/quality/sessions}
\end{itemize}

\subsection{Realtime APIs}

\begin{itemize}
\item \texttt{POST /v1/realtime/sessions}
\item \texttt{GET /v1/realtime/sessions/{sessionId}}
\item \texttt{POST /v1/realtime/sessions/{sessionId}/finalize}
\item \texttt{POST /v1/realtime/sessions/{sessionId}/confirm}
\item \texttt{POST /v1/realtime/sessions/{sessionId}/cancel}
\item \texttt{WSS /v1/realtime/sessions/{sessionId}/stream}
\end{itemize}

\subsection{Approval APIs}

\begin{itemize}
\item \texttt{GET /v1/approvals}
\item \texttt{POST /v1/approvals/{approvalId}/approve}
\item \texttt{POST /v1/approvals/{approvalId}/reject}
\end{itemize}

\subsection{Webhook Event Types}

\begin{itemize}
\item \texttt{voice.session.started}
\item \texttt{voice.session.completed}
\item \texttt{workflow.draft.ready}
\item \texttt{workflow.confirmed}
\item \texttt{workflow.approval.requested}
\item \texttt{workflow.approved}
\item \texttt{workflow.rejected}
\item \texttt{workflow.execution.succeeded}
\item \texttt{workflow.execution.failed}
\end{itemize}

% =========================================================
\section{Enterprise Dashboard}

The dashboard should include the following modules.

\subsection{Overview}

\begin{itemize}
\item Active sessions.
\item Completed workflows.
\item Failed workflows.
\item Average end-to-end latency.
\item Extraction accuracy.
\item Confirmation rate.
\item Approval backlog.
\item Integration-delivery health.
\end{itemize}

\subsection{Workflow Builder}

\begin{itemize}
\item Workflow details.
\item Visual JSON Schema builder.
\item Advanced JSON editor.
\item Trigger phrases.
\item Example utterances.
\item Business terminology.
\item Validation rules.
\item Execution policy.
\item Integration target.
\item Version history.
\end{itemize}

\subsection{Voice Playground}

The Voice Playground should allow an administrator to:

\begin{itemize}
\item Record or upload voice.
\item View partial transcripts.
\item Inspect workflow candidates.
\item Inspect extracted fields and confidence.
\item Test corrections.
\item Review missing-field questions.
\item View validation results.
\item Compare expected and actual JSON.
\item Save examples to the evaluation dataset.
\end{itemize}

\subsection{Quality Center}

The Quality Center should provide:

\begin{itemize}
\item Failed-session review.
\item Low-confidence sessions.
\item Frequently corrected fields.
\item Workflow confusion matrix.
\item Terminology gaps.
\item Model-version comparison.
\item Regression-evaluation reports.
\end{itemize}

\subsection{Integrations}

\begin{itemize}
\item Webhook destinations.
\item API connectors.
\item n8n templates.
\item Secret rotation.
\item Delivery logs.
\item Replay controls.
\item Test-event delivery.
\end{itemize}

% =========================================================
\section{Observability}

\subsection{Technical Metrics}

The platform should monitor:

\begin{itemize}
\item Active realtime connections.
\item Audio-stream duration.
\item Speech-service latency.
\item Partial transcript latency.
\item Workflow extraction latency.
\item Validation latency.
\item Model token usage.
\item Queue depth.
\item Webhook success rate.
\item Database and Redis performance.
\end{itemize}

\subsection{Product Quality Metrics}

\begin{itemize}
\item Workflow classification accuracy.
\item Field extraction precision.
\item Field extraction recall.
\item Correction frequency.
\item Missing-field frequency.
\item User confirmation rate.
\item User cancellation rate.
\item Manual-edit rate.
\item Successful execution rate.
\end{itemize}

\subsection{Tracing}

Each session should have a distributed trace ID propagated through:

\begin{itemize}
\item Client.
\item Realtime Gateway.
\item Speech Service.
\item Workflow Runtime.
\item Validation Service.
\item Dispatcher.
\item Webhook Worker.
\end{itemize}

% =========================================================
\section{Deployment Topology}

A production deployment may use:

\begin{itemize}
\item Kubernetes or another container orchestration platform.
\item Regional realtime gateways.
\item Dedicated speech-worker pools.
\item Autoscaled workflow-runtime workers.
\item Managed PostgreSQL with high availability.
\item Redis with replication or clustering.
\item Durable event streaming or queues.
\item S3-compatible object storage.
\item Centralized secrets management.
\item OpenTelemetry-compatible observability.
\end{itemize}

Enterprise deployment options may include:

\begin{itemize}
\item Multi-tenant managed cloud.
\item Dedicated tenant environment.
\item Private cloud.
\item Customer-controlled self-hosted deployment.
\item Hybrid edge and cloud deployment.
\end{itemize}

% =========================================================
\section{MVP Scope}

\subsection{Included in the MVP}

The first commercially testable version should include:

\begin{itemize}
\item Multi-tenant organization management.
\item User authentication and role-based access.
\item Workflow creation.
\item JSON Schema configuration.
\item Workflow examples and terminology.
\item Four initial workflow types:
purchase, sale, income, and expense.
\item Arabic and Egyptian Arabic streaming transcription.
\item WebSocket realtime sessions.
\item Incremental workflow drafts.
\item Missing-field detection.
\item Deterministic validation.
\item User confirmation.
\item Signed webhook delivery.
\item n8n-compatible webhook templates.
\item Session and delivery audit logs.
\item Voice testing playground.
\item Basic quality and latency dashboard.
\end{itemize}

\subsection{Excluded from the MVP}

The following should be postponed unless required by a paying pilot
customer:

\begin{itemize}
\item Full offline speech recognition.
\item Customer-specific model fine-tuning.
\item On-premise installation.
\item Advanced visual rule engine.
\item Multi-step workflow orchestration.
\item Marketplace of public workflow templates.
\item Automated financial posting.
\item Autonomous execution of high-risk actions.
\end{itemize}

% =========================================================
\section{Delivery Roadmap}

\begin{tabularx}{\textwidth}{L{28mm}Y Y}
\toprule
\textbf{Phase} & \textbf{Objectives} & \textbf{Main Deliverables} \\
\midrule
Phase 0:
Validation
&
Confirm customer problem and technical feasibility.
&
Pilot workflows, voice samples, latency benchmark, security assumptions. \\

Phase 1:
Foundation
&
Build tenant, IAM, workflow configuration, and schema versioning.
&
Management APIs, dashboard foundation, PostgreSQL model. \\

Phase 2:
Realtime Runtime
&
Build secure streaming voice sessions.
&
Realtime Gateway, speech integration, Redis session state. \\

Phase 3:
Structured Extraction
&
Convert partial transcripts into workflow drafts.
&
Workflow classification, field extraction, correction handling. \\

Phase 4:
Validation
&
Add deterministic business validation and entity resolution.
&
Validation engine, lookup connectors, missing-field handling. \\

Phase 5:
Execution
&
Safely deliver confirmed workflows.
&
Webhook dispatcher, retries, signatures, n8n templates. \\

Phase 6:
Enterprise Readiness
&
Improve auditability, quality, scale, and security.
&
Quality center, tracing, retention policies, load tests. \\
\bottomrule
\end{tabularx}

% =========================================================
\section{Team Responsibilities}

\begin{tabularx}{\textwidth}{L{42mm}Y}
\toprule
\textbf{Team or Role} & \textbf{Responsibilities} \\
\midrule
Product Management
& Customer discovery, workflow priorities, pricing, acceptance criteria. \\

Backend Engineering
& APIs, tenancy, workflow runtime, validation, integrations, persistence. \\

Realtime Engineering
& WebSocket/WebRTC gateway, audio streaming, session recovery, latency. \\

AI and ML Engineering
& Speech models, extraction strategy, evaluation, model routing. \\

Frontend Engineering
& Enterprise dashboard, schema builder, voice playground, session UI. \\

Security Engineering
& Threat modeling, IAM, tenant isolation, keys, audit, compliance. \\

DevOps and SRE
& Deployment, autoscaling, observability, reliability, incident response. \\

QA Engineering
& Functional tests, realtime tests, regression datasets, load tests. \\

Domain Experts
& Business rules, terminology, workflow examples, validation review. \\
\bottomrule
\end{tabularx}

% =========================================================
\section{Testing Strategy}

\subsection{Functional Testing}

\begin{itemize}
\item Workflow configuration.
\item Schema validation.
\item Workflow version publishing.
\item Permission enforcement.
\item Session creation and completion.
\item Confirmation and approval.
\item Webhook delivery and replay.
\end{itemize}

\subsection{Voice and AI Evaluation}

The evaluation dataset should include:

\begin{itemize}
\item Egyptian Arabic accents.
\item Background noise.
\item Slow and fast speakers.
\item Corrections.
\item Multiple line items.
\item Similar customer and product names.
\item Mixed Arabic and English terminology.
\item Spoken numbers and monetary amounts.
\item Incomplete statements.
\item Ambiguous workflows.
\end{itemize}

\subsection{Performance Testing}

Performance testing should cover:

\begin{itemize}
\item Concurrent WebSocket sessions.
\item Long-running audio streams.
\item Reconnection storms.
\item Speech-worker saturation.
\item Model-provider throttling.
\item Redis failover.
\item Queue backlog.
\item Webhook destination failures.
\end{itemize}

\subsection{Security Testing}

\begin{itemize}
\item Cross-tenant access attempts.
\item Token theft and replay.
\item Prompt-injection attempts.
\item Webhook forgery.
\item Secret leakage.
\item Malicious audio and payload sizes.
\item Permission bypass attempts.
\item Sensitive-data logging.
\end{itemize}

% =========================================================
\section{MVP Acceptance Criteria}

The MVP may be considered ready for a controlled customer pilot when:

\begin{enumerate}
\item An enterprise administrator can define and publish a workflow
without modifying backend code.

\item The administrator can provide a JSON Schema, examples,
terminology, and validation rules.

\item An authenticated employee can start a realtime voice session.

\item Partial transcript updates are displayed while the employee speaks.

\item The system incrementally identifies one of the configured workflows.

\item The system produces a structured draft matching the published schema.

\item Spoken corrections update the existing draft instead of creating
duplicate actions.

\item Missing required fields are identified and requested.

\item Deterministic validation prevents invalid final actions.

\item High-risk workflows cannot execute without their configured
confirmation or approval.

\item Confirmed workflows are delivered through signed, idempotent
webhooks.

\item Failed webhook deliveries are retried and visible in the dashboard.

\item Every final action is traceable through an immutable audit record.

\item Tenant data cannot be accessed by another tenant.

\item The realtime system satisfies the agreed pilot latency and
concurrency targets.
\end{enumerate}

% =========================================================
\section{Primary Risks and Mitigations}

\begin{longtable}{L{45mm}L{61mm}L{55mm}}
\toprule
\textbf{Risk} & \textbf{Impact} & \textbf{Mitigation} \\
\midrule
\endhead

Poor dialect recognition
& Incorrect transcript and extraction
& Tenant terminology, evaluation datasets, model routing, confirmations \\

High model latency
& Unresponsive experience
& Fast extraction layer, semantic checkpoints, caching, regional workers \\

Incorrect action execution
& Financial or operational damage
& Draft-only processing, deterministic validation, confirmation, approval \\

Cross-tenant leakage
& Critical security breach
& Strict tenant context, row-level enforcement, isolated keys and logs \\

Prompt injection
& Invalid or unauthorized model behavior
& Constrained outputs, limited context, independent authorization \\

External integration failure
& Lost or duplicate actions
& Idempotency, retries, dead-letter queues, delivery logs \\

Cost growth
& Unsustainable margins
& Model routing, audio limits, batching, per-tenant quotas, usage pricing \\

Workflow complexity
& Difficult enterprise onboarding
& Visual builder, templates, guided testing, professional services \\

Noisy environments
& Poor speech quality
& Voice activity detection, noise handling, client microphone guidance \\

User distrust
& Low adoption
& Live draft preview, explainable corrections, final confirmation \\
\bottomrule
\end{longtable}

% =========================================================
\section{Commercial Model}

A hybrid commercial model is recommended.

\subsection{Platform Subscription}

The subscription may be based on:

\begin{itemize}
\item Number of active users.
\item Number of voice minutes.
\item Number of completed workflows.
\item Number of published workflow definitions.
\item Data-retention requirements.
\item Deployment model.
\end{itemize}

\subsection{Usage Charges}

Usage charges may cover:

\begin{itemize}
\item Streaming speech minutes.
\item Structured-model processing.
\item Webhook events.
\item Stored audio.
\item Premium entity-resolution connectors.
\end{itemize}

\subsection{Enterprise Services}

Additional revenue may come from:

\begin{itemize}
\item Workflow design.
\item Integration implementation.
\item Custom deployment.
\item Security review.
\item Domain-specific terminology preparation.
\item Evaluation-dataset preparation.
\item Service-level agreements.
\end{itemize}

% =========================================================
\section{Recommended Initial Pilot}

The first pilot should focus on one operational customer and four
high-frequency workflows:

\begin{enumerate}
\item Purchase
\item Sale
\item Income or receipt
\item Expense or payment
\end{enumerate}

The pilot should avoid direct irreversible accounting posting.

The initial integration should create a validated transaction draft or
send a signed webhook to the customer's backend. The customer's system
remains the final system of record.

This approach reduces risk while proving:

\begin{itemize}
\item Real-time usability.
\item Arabic speech quality.
\item Workflow configurability.
\item Integration value.
\item Operational time savings.
\item Customer willingness to pay.
\end{itemize}

% =========================================================
\section{Final Recommendation}

The project should be developed as a secure, multi-tenant,
configuration-driven runtime rather than as a collection of hard-coded
voice commands.

The architecture should preserve a strict separation between:

\begin{itemize}
\item Speech recognition.
\item AI-based extraction.
\item Deterministic validation.
\item Authorization.
\item Confirmation and approval.
\item Final execution.
\end{itemize}

The first product milestone should prove one complete vertical flow:

\begin{center}
\textbf{
Tenant Workflow Definition
$\rightarrow$
Realtime Voice Session
$\rightarrow$
Structured Draft
$\rightarrow$
Validation
$\rightarrow$
User Confirmation
$\rightarrow$
Signed Webhook
}
\end{center}

Once this flow is stable, measurable, and secure, the platform can expand
into additional industries, languages, deployment models, and workflow
types.

% =========================================================
\appendix

\section{Decision Summary}

\begin{tabularx}{\textwidth}{L{53mm}Y}
\toprule
\textbf{Decision} & \textbf{Selected Direction} \\
\midrule
Primary product category
& Enterprise realtime voice workflow infrastructure \\

Customization model
& Tenant-defined schemas, examples, terminology, and rules \\

Realtime transport
& WebSocket initially; WebRTC where audio requirements justify it \\

Execution model
& Streaming draft, final confirmed execution \\

State model
& Redis for hot state, PostgreSQL for durable state \\

Integration model
& Signed webhooks, REST APIs, queues, and n8n \\

AI role
& Classification and extraction, not authorization or final validation \\

MVP workflows
& Purchase, sale, income, and expense \\

Initial delivery risk
& Create validated drafts or webhooks, not direct financial posting \\

Security model
& Tenant isolation, RBAC, audit, encryption, signed delivery \\
\bottomrule
\end{tabularx}

\section{Glossary}

\begin{description}[style=nextline,leftmargin=35mm]
\item[Tenant]
A customer organization isolated within the platform.

\item[Workflow Definition]
A versioned configuration describing a business operation and its
structured data.

\item[Workflow Draft]
The mutable structured result produced while the employee is speaking.

\item[Voice Session]
A realtime interaction containing audio, transcripts, drafts,
validations, and execution state.

\item[Entity Resolution]
Mapping spoken names to valid customer, supplier, product, account, or
other business identifiers.

\item[Execution Policy]
Rules controlling confirmation, authentication, approval, and final
action execution.

\item[Semantic Checkpoint]
A meaningful point at which the reasoning model updates the structured
workflow draft.

\item[System of Record]
The authoritative customer application that owns the final business
transaction.
\end{description}

\end{document}
